% a01_H0rel part 6x3, continued: the count over the paid elements (input from a01_p6x3.tex after lem:a01:PP)
\begin{lemma}[the count over the paid elements]\label{lem:a01:PP5}
Let the boxes $\mathcal Q_i$, the relation $\sim$, the arcs of $q_i$ to $2q_i-1$ large blocks, the set $\mathcal C_i$ and the family $F_i$ be those of \cref{lem:a01:PP}, with the charge point $x_e$ of an element $e$ as in its proof (every point of $e$ lies within $\dist\le L-1$ of $x_e$). Let $F_i^{\mathrm{paid}}\subset F_i$, and suppose that every $e\in F_i\setminus F_i^{\mathrm{paid}}$ has a \emph{witness} $w(e)\in F_i^{\mathrm{paid}}$ with $\dist(x_e,x_{w(e)})\le r_i+s_b$. Then
\[
|\mathcal C_i|\ \le\ 5^d\,|F_i^{\mathrm{paid}}|.
\]
In D14 the unpaid elements are the near $R_s$ elements, the blocks $y\in R_s^{(i)}$ within $\ell^\infty$-distance $R_0$ of the face layer, and $F_i^{\mathrm{paid}}$ consists of Ba{\l}aban's six families, $R_s$ restricted to its far elements, and the face family; the witness of D7(b) (\cref{lem:a01:F541}(e)) has a point within distance $dR_0+1$ of a point of $e$, so $\dist(x_e,x_{w(e)})\le dR_0+1+2(L-1)\le s_b+L-2<r_i+s_b$ by $s_b\ge dR_0+L+1$ (D7(c)) and $r_i\ge R>L$, and the hypothesis holds.
\end{lemma}
\begin{proof}
Let $Q\in\mathcal C_i$. The proof of \cref{lem:a01:PP} gives $z'\in Q$ and $e\in F_i$ with $\dist(z',x_e)\le r_i+s_b$. Put $e'\defeq e$ if $e\in F_i^{\mathrm{paid}}$ and $e'\defeq w(e)$ otherwise; in both cases $e'\in F_i^{\mathrm{paid}}$ and $\dist(z',x_{e'})\le2(r_i+s_b)\le2q_is_b$, since $q_is_b\ge r_i+s_b$. In each coordinate the arc indices of $Q=Q(z')$ and of $Q(x_{e'})$ differ by at most $2$: if they differed by $3$ or more, two whole arcs, of at least $q_is_b$ points each, would lie strictly between the two coordinates, whose distance would then be $\ge2q_is_b+1$. A box has at most $5^d$ boxes whose arc indices differ from its own by at most $2$ in every coordinate (fewer on a torus with fewer than five arcs in some coordinate). So $Q\mapsto e'$ maps $\mathcal C_i$ into $F_i^{\mathrm{paid}}$ with fibres of at most $5^d$ boxes, which is the claim. For the application: a near $R_s$ element $e$ at the block $y$ has $\dist(y,x_e)\le L-1$; \cref{lem:a01:F541}(e) gives, on a block path of length $\le dR_0$ from $y$ to the face layer, a bond of the gradient family or a block of the face family, that is an element $w(e)$ of a paid family one of whose points lies within distance $dR_0+1$ of $y$ and within $L-1$ of $x_{w(e)}$; the triangle inequality gives $\dist(x_e,x_{w(e)})\le dR_0+2L-1$, and $dR_0+2L-1\le s_b+L-2$ is $s_b\ge dR_0+L+1$.
\end{proof}
